% 附录A：算法分类速查表与学习路线图
\chapter{算法分类速查表与学习路线图}

本附录汇总全书涉及的主要算法，提供分类速查和学习建议。

% =============================================
\section*{A.1 算法总表}

\begin{longtable}{p{3.5cm}p{4cm}p{1.2cm}p{4.5cm}}
\toprule
\textbf{算法类别} & \textbf{代表算法} & \textbf{层次} & \textbf{核心特点} \\
\midrule
\endfirsthead
\toprule
\textbf{算法类别} & \textbf{代表算法} & \textbf{层次} & \textbf{核心特点} \\
\midrule
\endhead
\bottomrule
\endfoot

\multicolumn{4}{l}{\textbf{第一篇：经典规划方法}} \\
\midrule
无信息搜索 & BFS, DFS, IDS & \textcolor{ugcolor}{本} & 不使用启发信息 \\
启发式搜索 & A*, IDA*, GBFS & \textcolor{ugcolor}{本} & 启发函数引导搜索 \\
加权A* & WA* & \textcolor{gradcolor}{研} & 牺牲最优性换速度 \\
松弛启发式 & $h^{\max}$, $h^{\mathrm{add}}$, $h^{\mathrm{FF}}$ & \textcolor{ugcolor}{本} & 删除松弛估计 \\
规划图 & GraphPlan & \textcolor{ugcolor}{本} & 互斥分析+解提取 \\
抽象启发式 & PDB, Merge\&Shrink & \textcolor{gradcolor}{研} & 问题投影和抽象 \\
地标启发式 & LM-Cut, 地标计数 & \textcolor{gradcolor}{研} & 必经之路分析 \\
SAT规划 & SATplan, Madagascar & \textcolor{ugcolor}{本} & 转化为SAT求解 \\
CSP/SMT规划 & CSP编码, SMT规划 & \textcolor{ugcolor}{本}/\textcolor{gradcolor}{研} & 约束满足求解 \\
HTN规划 & SHOP/SHOP2, PANDA & \textcolor{ugcolor}{本} & 层次任务分解 \\
偏序规划 & POP & \textcolor{ugcolor}{本} & 最小承诺策略 \\
时态规划 & TFD, OPTIC & \textcolor{gradcolor}{研} & 处理持续性动作 \\
MDP/POMDP & 值迭代, PBVI & \textcolor{gradcolor}{研} & 概率不确定性 \\
FF规划器 & FF (EHC+$h^{\mathrm{FF}}$) & \textcolor{ugcolor}{本} & 松弛启发+爬山 \\
Fast Downward & FD (SAS$^+$) & \textcolor{ugcolor}{本} & 模块化框架 \\
LAMA & 地标+多启发式 & \textcolor{ugcolor}{本} & 满足规划冠军 \\
Scorpion/SymK & 饱和代价/符号搜索 & \textcolor{gradcolor}{研} & 最优规划前沿 \\
\midrule

\multicolumn{4}{l}{\textbf{第二篇：元启发式与优化方法}} \\
\midrule
遗传算法 & GA & \textcolor{ugcolor}{本} & 选择交叉变异 \\
模拟退火 & SA & \textcolor{ugcolor}{本} & 概率接受+冷却 \\
差分进化 & DE & \textcolor{gradcolor}{研} & 差异向量驱动 \\
禁忌搜索 & TS & \textcolor{gradcolor}{研} & 禁忌表+特赦 \\
粒子群优化 & PSO & \textcolor{gradcolor}{研} & 群体信息共享 \\
蚁群优化 & ACO, MMAS & \textcolor{gradcolor}{研} & 信息素引导 \\
新型元启发式 & GWO, WOA, HHO & \textcolor{selfcolor}{拓} & 自然现象启发 \\
\midrule

\multicolumn{4}{l}{\textbf{第三篇：强化学习与规划}} \\
\midrule
Dyna架构 & Dyna-Q & \textcolor{ugcolor}{本} & 学习+规划统一 \\
蒙特卡洛树搜索 & MCTS, UCT & \textcolor{ugcolor}{本} & 采样+模拟评估 \\
AlphaGo/Zero & MCTS+深度网络 & \textcolor{gradcolor}{研} & 自我对弈+搜索 \\
MuZero & 学习的世界模型 & \textcolor{gradcolor}{研} & 隐空间规划 \\
Dreamer系列 & DreamerV1--V3 & \textcolor{gradcolor}{研} & 想象中的规划 \\
DQN系列 & DQN, Rainbow & \textcolor{ugcolor}{本} & 深度Q网络 \\
策略梯度 & REINFORCE, A2C & \textcolor{ugcolor}{本} & 直接优化策略 \\
PPO & PPO & \textcolor{gradcolor}{研} & 裁剪策略更新 \\
SAC & SAC & \textcolor{gradcolor}{研} & 最大熵RL \\
层次RL & Options, HIRO & \textcolor{gradcolor}{研} & 时间抽象 \\
Decision Transformer & DT & \textcolor{gradcolor}{研} & 序列建模=RL \\
\midrule

\multicolumn{4}{l}{\textbf{第四篇：深度学习与规划}} \\
\midrule
图神经网络规划 & ASNets & \textcolor{gradcolor}{研} & 图结构泛化 \\
Transformer规划 & Plansformer & \textcolor{gradcolor}{研} & 端到端生成 \\
扩散模型规划 & Diffuser, Diff.\ Policy & \textcolor{selfcolor}{拓} & 轨迹去噪生成 \\
神经符号规划 & NLM, NSRT & \textcolor{selfcolor}{拓} & 感知+推理融合 \\
\midrule

\multicolumn{4}{l}{\textbf{第五篇：大语言模型与规划}} \\
\midrule
LLM直接规划 & Zero/Few-shot规划 & \textcolor{gradcolor}{研} & 语言知识驱动 \\
LLM+规划器 & LLM+P & \textcolor{gradcolor}{研} & 混合求解 \\
思维链/树/图 & CoT, ToT, GoT & \textcolor{gradcolor}{研} & 结构化推理 \\
具身LLM规划 & SayCan, Code as Policies & \textcolor{selfcolor}{拓} & 机器人+LLM \\
语言智能体 & ReAct, Reflexion & \textcolor{gradcolor}{研} & 推理+行动+反思 \\
多智能体LLM & AutoGen, MetaGPT & \textcolor{selfcolor}{拓} & 角色协作 \\
自主探索 & Voyager & \textcolor{selfcolor}{拓} & 终身学习 \\
\midrule

\multicolumn{4}{l}{\textbf{第六篇：多智能体规划}} \\
\midrule
MAPF & CBS, ICTS & \textcolor{ugcolor}{本}/\textcolor{gradcolor}{研} & 无冲突路径 \\
LaCAM & LaCAM/LaCAM* & \textcolor{gradcolor}{研} & 大规模MAPF \\
博弈论方法 & Nash, Stackelberg & \textcolor{gradcolor}{研} & 利益冲突建模 \\
MARL & QMIX, MAPPO & \textcolor{gradcolor}{研} & 集中训练分布执行 \\
分布式规划 & MA-STRIPS & \textcolor{selfcolor}{拓} & 隐私保护协调 \\
\end{longtable}

% =============================================
\section*{A.2 推荐学习路线}

\subsection*{本科生核心路线（建议32--48学时）}

\begin{enumerate}
  \item \textbf{基础（8--12学时）：}第1章（搜索基础）$\rightarrow$ 第2章1--2节（松弛启发式、规划图）
  \item \textbf{经典算法（8--12学时）：}第3章1--4节（SAT、HTN、POP）$\rightarrow$ 第4章1--3节（FF、FD、LAMA）
  \item \textbf{优化方法（4--6学时）：}第5章1--2节（GA、SA）
  \item \textbf{强化学习（8--12学时）：}第6章1--2节（Dyna、MCTS）$\rightarrow$ 第7章1--2节（DQN、策略梯度）
  \item \textbf{多智能体（4--6学时）：}第11章1节（MAPF、CBS）
\end{enumerate}

\subsection*{研究生扩展路线（在本科基础上增加16--32学时）}

\begin{enumerate}
  \item \textbf{高级启发式（4--6学时）：}第2章3--4节（PDB、LM-Cut）
  \item \textbf{时态与不确定性（4--6学时）：}第3章5--6节
  \item \textbf{高级RL（4--8学时）：}第6章3--4节（AlphaZero、Dreamer）$\rightarrow$ 第7章3--4节（层次RL、DT）
  \item \textbf{深度学习规划（4--6学时）：}第8章1--2节（GNN、Transformer）
  \item \textbf{LLM规划（4--8学时）：}第9章1--3节 $\rightarrow$ 第10章1--2节
  \item \textbf{多智能体进阶（4--6学时）：}第11章2--3节（博弈论、MARL）
\end{enumerate}

\subsection*{自学拓展专题}

以下内容适合感兴趣的读者根据研究方向选择阅读：
\begin{itemize}
  \item 扩散模型与规划（第8章3节）
  \item 神经符号规划（第8章4节）
  \item 具身智能LLM规划（第9章4节）
  \item 语言智能体前沿（第10章3--4节）
  \item 分布式规划（第11章4节）
  \item 新型元启发式综述（第5章4节）
\end{itemize}

% =============================================
\section*{A.3 算法演化脉络图}

任务规划算法的发展可以概括为以下几条主线：

\begin{enumerate}
  \item \textbf{符号推理主线：}STRIPS $\rightarrow$ GraphPlan $\rightarrow$ SAT规划 $\rightarrow$ 启发式搜索（FF/FD） $\rightarrow$ 地标/抽象启发式
  \item \textbf{层次分解主线：}HTN $\rightarrow$ SHOP/SHOP2 $\rightarrow$ PANDA $\rightarrow$ 时态HTN
  \item \textbf{优化求解主线：}SA/GA $\rightarrow$ PSO/ACO $\rightarrow$ 混合元启发式
  \item \textbf{学习规划主线：}Dyna $\rightarrow$ DQN $\rightarrow$ AlphaGo/MuZero $\rightarrow$ Dreamer $\rightarrow$ Decision Transformer
  \item \textbf{深度生成主线：}VAE/GAN $\rightarrow$ Transformer $\rightarrow$ 扩散模型 $\rightarrow$ 轨迹生成规划
  \item \textbf{语言智能主线：}GPT $\rightarrow$ CoT/ToT $\rightarrow$ ReAct/Reflexion $\rightarrow$ Voyager $\rightarrow$ 多智能体LLM
\end{enumerate}

这些主线并非孤立发展，而是在不断交叉融合。例如，MuZero结合了学习和搜索，LLM+P结合了语言模型和符号规划，ASNets结合了深度学习和经典规划。未来的发展趋势是各主线的进一步融合，形成更强大的统一规划框架。

% =============================================
\section*{A.4 国内推荐参考文献}

\subsection*{教材与专著}
\begin{enumerate}
  \item 蔡自兴, 刘丽珏, 蔡竞峰.《人工智能及其应用》（第7版）. 清华大学出版社, 2024. ——国内使用最广泛的AI教材之一，含智能规划章节。
  \item 吴飞, 潘云鹤.《人工智能引论》. 高等教育出版社, 2024. ——``101计划''核心教材，体现国内AI教育前沿。
  \item Ghallab等著, 殷建平等译.《自动规划：理论和实践》. 清华大学出版社, 2008. ——国际经典教材中文版，系统讲述自动规划理论。
  \item 段海滨.《蚁群算法原理及其应用》. 科学出版社, 2005. ——蚁群算法的理论与工程应用。
  \item 焦李成等.《免疫优化计算、学习与识别》. 科学出版社, 2006. ——免疫计算与优化理论。
  \item 焦李成等.《深度学习、优化与识别》. 清华大学出版社, 2017. ——深度学习与优化方法。
  \item 陈杰, 辛斌.《智能优化的探索--开发权衡理论与方法》. 科学出版社, 2013. ——元启发式核心理论。
  \item 陈杰, 方浩, 辛斌.《多智能体系统的协同群集运动控制》. 科学出版社, 2017. ——多智能体协同控制。
  \item 向峥嵘等.《多智能体协同控制及应用》. 化学工业出版社, 2023. ——多智能体系统理论与应用。
  \item 余伶俐, 周开军, 陈白帆.《智能驾驶技术：路径规划与导航控制》. 机械工业出版社, 2020. ——自动驾驶中的规划技术。
  \item 赵鑫等.《大语言模型》. 电子工业出版社, 2024. ——大语言模型技术全面参考。
\end{enumerate}

\subsection*{重要综述论文}
\begin{enumerate}
  \item 高阳, 陈世福, 陆鑫. 强化学习研究综述. 自动化学报, 2004. ——中文RL经典入门综述。
  \item 刘全等. 深度强化学习综述. 计算机学报, 2018. ——深度RL方法体系梳理。
  \item 赵冬斌等. 深度强化学习综述：兼论计算机围棋的发展. 控制理论与应用, 2016. ——AlphaGo技术解析。
  \item 梁星星, 冯旸赫等. 多Agent深度强化学习综述. 自动化学报, 2020. ——MARL方法全面综述。
  \item 杜永浩, 邢立宁, 蔡昭权. 无人飞行器集群智能调度技术综述. 自动化学报, 2020. ——无人机调度方法体系。
  \item 贾高伟, 王建峰. 无人机集群任务规划方法研究综述. 系统工程与电子技术, 2021. ——无人机集群规划综述。
  \item 邢立宁, 陈英武. 任务规划系统研究及发展. 火力与指挥控制, 2006. ——任务规划系统发展脉络。
  \item 秦龙等. 基于大语言模型的复杂任务自主规划处理框架. 自动化学报, 2024. ——LLM任务规划前沿。
  \item 张宏宏等. 有人/无人机协同作战：概念、技术与挑战. 航空学报, 2024. ——有人/无人协同综述。
  \item 王建峰等. 多无人机协同任务规划方法研究综述. 系统工程与电子技术, 2024. ——多无人机协同最新进展。
  \item 王雪松等. 安全强化学习综述. 自动化学报, 2023. ——安全RL约束方法。
\end{enumerate}
